% !TEX program = tectonic
\documentclass[10.5pt]{article}
\usepackage{newtxtext,newtxmath}
\usepackage[a4paper,margin=0.86in]{geometry}
\usepackage{amsmath}
\usepackage{booktabs,longtable,array}
\usepackage{xcolor,enumitem,titlesec}
\usepackage[hidelinks,colorlinks=true,linkcolor=blue!55!black,urlcolor=teal!60!black]{hyperref}

\definecolor{g5}{HTML}{C65335}
\definecolor{teal}{HTML}{267B72}
\newcommand{\rv}{\mathbf r}
\newcommand{\cv}{\mathbf c}
\newcommand{\Fv}{\mathbf F}
\newcommand{\ev}{\widehat{\mathbf e}}
\newcommand{\keylabel}{\par\smallskip\noindent\textcolor{g5}{\textbf{KEY EQUATION}}\par\vspace{-5pt}}
\newcommand{\unit}[1]{\,\mathrm{#1}}
\setlist[itemize]{leftmargin=1.35em,itemsep=2pt,topsep=3pt}
\titleformat{\section}{\Large\bfseries\color{black!85}}{\thesection}{0.55em}{}
\titleformat{\subsection}{\large\bfseries\color{teal}}{\thesubsection}{0.55em}{}
\allowdisplaybreaks

\title{\vspace{-1.2cm}\textbf{The Generation-5 Tumour-Organoid Model}\\[3pt]
\large Theory, Governing Equations, Lineage from G4, and Ablation Logic\\[6pt]
\normalsize\textnormal{Motor--clutch collagen-remodelling project}}
\author{}
\date{2026-09-07}

\begin{document}
\maketitle
\vspace{-1cm}

\begin{abstract}
Generation 5 (G5) scales the validated single-cell motor--clutch and bead--spring collagen
framework to a multicellular tumour organoid. The clean conceptual change is one cell to many
mechanically interacting cell disks. This requires organoid geometry, cell--cell repulsion and
adhesion, repeated cell--ECM contact bookkeeping, and collective migration metrics. Tensile
strain stiffening, leader-cell heterogeneity, and crosslink topology plasticity are separate
hypotheses that should be activated by ablation rather than hidden in the baseline. The model is
two-dimensional and coarse-grained; its parameters are mechanism-first until calibrated against
organoid experiments.
\end{abstract}

\keylabel
\begin{equation}
 \zeta_b\dot{\rv}_i=
 \Fv_i^{\mathrm{stretch}}+\Fv_i^{\mathrm{bend}}+\Fv_i^{\mathrm{xlink}}
 +\sum_a\Fv_{ia}^{\mathrm{rep}}+
 \sum_a\Fv_{ia}^{\mathrm{cell\to ECM}}.
 \label{eq:backbone}
\end{equation}
The collagen equation remains the common backbone. G5 changes the number of cells and the
coupling terms, not the unit system.

\tableofcontents

\section{Lineage: what is inherited and what is new}

G5 inherits the G2/G4 bead-chain geometry, axial stretch, discrete bending, compression
softening, outer-boundary anchoring, permanent elastic crosslinks, local material-point contact,
and optional Bell motor--clutch kinetics. The genuinely required additions are multiple rigid
cell disks, cell--cell mechanics, per-cell contact ownership, and collective outcomes.

\begin{longtable}{>{\raggedright\arraybackslash}p{0.16\textwidth}p{0.28\textwidth}p{0.45\textwidth}}
\toprule
Stage & New mechanism & Primary question \\
\midrule
G5A & $N$ cell disks plus cohesion & Is the organoid stable before cell--ECM pulling? \\
G5B & Many cells pull one shared ECM & Does collective contraction increase radial collagen
alignment? \\
G5C & Tensile strain stiffening & Does nonlinear recruitment extend the deformation halo? \\
G5D & Reaction-driven cell motion & Does cell--cell adhesion select collective versus
single-cell escape? \\
G5E & Crosslink rupture and re-welding & Does topology change preserve remodelling after
unloading? \\
\bottomrule
\end{longtable}

The clean G5 baseline is multicellular disks, cell--cell adhesion, inherited elastic collagen,
and inherited clutch mechanics. Strain stiffening, leader cells, and crosslink plasticity are
explicit switches. This separation keeps G4 collagen calibration conceptually distinct from the
multicellular consequences introduced in G5.

\section{Inherited collagen mechanics}

Every collagen fibre is a connected sequence of beads. Visible lines join consecutive beads,
and crosslinks join material points on different fibres. The outer boundary is anchored so local
remodelling is not confused with rigid translation of the whole network.

For non-fixed beads,
\begin{equation}
 \rv_i(t+\Delta t)=\rv_i(t)+\frac{\Delta t}{\zeta_b}\Fv_i^{\mathrm{total}}.
 \label{eq:euler}
\end{equation}
There is no inertia and no standard-linear-solid element. Drag controls transient speed; elastic
terms store energy.

For an axial bond,
\begin{align}
 e_{ij}&=\ell_{ij}-\ell_{ij}^0,
 &\Fv_{i\leftarrow j}^{\mathrm{stretch}}&=k_{ij}e_{ij}\ev_{ij},\\
 k_{ij}&=
 \begin{cases}
 EA/\ell_{ij}^0,&e_{ij}\ge0,\\
 \rho EA/\ell_{ij}^0,&e_{ij}<0,
 \end{cases}
 &\ev_{ij}&=(\rv_j-\rv_i)/\ell_{ij}.
 \label{eq:stretch}
\end{align}
Compression remains soft ($\rho<1$), representing microbuckling rather than an equally stiff
compression spring.

For bending,
\begin{equation}
 \Delta\boldsymbol\kappa_i=
 (\rv_{i-1}-2\rv_i+\rv_{i+1})-\boldsymbol\kappa_i^0,
 \qquad
 (\Fv_{i-1}^{b},\Fv_i^{b},\Fv_{i+1}^{b})
 =k_b(-\Delta\boldsymbol\kappa_i,2\Delta\boldsymbol\kappa_i,-\Delta\boldsymbol\kappa_i).
 \label{eq:bending}
\end{equation}

A freely hinged crosslink joining two material points obeys
\begin{align}
 \mathbf p_a&=(1-\alpha_a)\rv_i+\alpha_a\rv_j,\\
 \Fv_x&=k_x\big[(\mathbf p_b-\mathbf p_a)-\mathbf d_{x,0}\big].
 \label{eq:xlink}
\end{align}
The link transmits separation force without prescribing the crossing angle.

\section{G5A: multicellular scaffold and resting balance}

Cell centres are placed on a hexagonal lattice with rest pitch
$d_0=18\unit{\mu m}$. Each rigid disk has radius $R_c=9\unit{\mu m}$, and centres inside the
organoid target radius are retained. Collagen is excluded from the union of cell disks plus a
clearance region.

Let $d=\lVert\cv_b-\cv_a\rVert$ and
$\ev_{ab}=(\cv_b-\cv_a)/d$. The pairwise force on cell $a$ from cell $b$ is
\keylabel
\begin{equation}
 \Fv_{a\leftarrow b}^{cc}=
 \begin{cases}
 -k_{\mathrm{rep}}(d_0-d)\ev_{ab},&d<d_0,\\
 k_{\mathrm{adh}}s_{ab}(d-d_0)\ev_{ab},&d_0\le d<d_0+r_{\mathrm{adh}},\\
 \mathbf 0,&d\ge d_0+r_{\mathrm{adh}},
 \end{cases}
 \qquad s_{ab}=\min(s_a,s_b).
 \label{eq:cellcell}
\end{equation}
Repulsion prevents overlap and finite-range adhesion supplies coarse tissue cohesion. The
strength $k_{\mathrm{adh}}$ is an effective organoid-scale parameter, not a direct cadherin
single-bond stiffness. At $d=d_0$, the pair force vanishes.

Cell--collagen exclusion generalizes the G4 penalty to all cells:
\begin{equation}
 \delta_{ia}=\max\big(0,R_c+c-\lVert\rv_i-\cv_a\rVert\big),\qquad
 \Fv_i^{\mathrm{rep}}=\sum_a k_{\mathrm{rep,ECM}}\delta_{ia}
 \frac{\rv_i-\cv_a}{\lVert\rv_i-\cv_a\rVert}.
 \label{eq:multirep}
\end{equation}
The G5A gate is a stable non-overlapping organoid that neither collapses nor disperses before
active traction is introduced.

\section{G5B: fixed organoid pulling a shared collagen network}

For cell $a$ and fibre material point $\mathbf p_{af}$, contact is eligible only inside a finite
surface shell. Gaussian weights are normalized over the eligible set $\mathcal C(a)$:
\keylabel
\begin{align}
 s_{af}&=\lVert\mathbf p_{af}-\cv_a\rVert-R_c,
 &0&\le s_{af}\le w_c,\\
 w_{af}&=\frac{\exp(-s_{af}^2/\sigma_c^2)}
 {\displaystyle\sum_{g\in\mathcal C(a)}\exp(-s_{ag}^2/\sigma_c^2)}.
 \label{eq:contact}
\end{align}
The inward pull from cell $a$ is
\begin{equation}
 \Fv_{af}^{\mathrm{cell\to ECM}}
 =F_a w_{af}\frac{\cv_a-\mathbf p_{af}}
 {\lVert\cv_a-\mathbf p_{af}\rVert}.
 \label{eq:pull}
\end{equation}
Linear shape functions distribute this point force to the two host beads. A non-contact fibre
receives no second Gaussian pull; it can move only through crosslink force transmission and its
own stretching and bending. Removing crosslinks is therefore the causal negative control.

Radial alignment is measured on each fibre edge by
\keylabel
\begin{equation}
 S_{r,e}=2(\mathbf t_e\cdot\ev_{r,e})^2-1,
 \qquad
 S_r(q)=\frac{1}{N_q}\sum_{e\in q}S_{r,e}.
 \label{eq:radial}
\end{equation}
$S_r=+1$ is radial, $S_r=0$ isotropic, and $S_r=-1$ tangential. Reporting distance shells
separates near-field reorientation from a whole-network average.

\section{G5C: tensile strain-stiffening ablation}

The optional nonlinear constitutive law changes only tensile bonds:
\keylabel
\begin{equation}
 \varepsilon=\frac{e}{\ell_0},\qquad
 k_t(\varepsilon)=k_{t,0}
 \min\!\left(C_{\max},\exp\!\frac{\varepsilon}{\varepsilon_{\mathrm{ref}}}\right)
 \quad(\varepsilon\ge0),
 \qquad
 k(\varepsilon<0)=\rho k_{t,0}.
 \label{eq:stiffening}
\end{equation}
Current defaults use $\varepsilon_{\mathrm{ref}}=0.06$ and $C_{\max}=12$ as exploratory
guards. A taut bond becomes stiffer and can recruit a longer load-bearing tract. The clean
comparison switches strain stiffening off and on while holding network, pull, cells, and seed
fixed. If most strains remain only one to five percent, the nonlinear branch is only weakly
engaged; a small effect is then a statement about the current loading regime, not a rejection of
collagen nonlinearity.

\section{Motor--clutch kinetics in a multicellular setting}

When clutch dynamics are enabled, every cell is divided into angular contact sectors and every
occupied sector contains an effective clutch bundle. G5 reuses G4 kinetics but repeats them over
many cells and preserves site ownership for reaction-force bookkeeping.

Motor loading follows
\keylabel
\begin{equation}
 v_{\mathrm{act}}=v_0\max\!\left(0,1-\frac{F_{\mathrm{site}}}{F_{\mathrm{stall}}}\right),
 \qquad
 \dot x=\max(0,v_{\mathrm{act}}-v_{\mathrm{sub}}),
 \qquad f_j=k_cx_j.
 \label{eq:motor}
\end{equation}

Bell slip-bond rupture and rebinding are
\begin{equation}
 k_{\mathrm{off}}(f)=k_{\mathrm{off}}^0\exp\!\left(\frac{f}{F_b}\right),
 \quad P_{\mathrm{off}}=1-\exp[-k_{\mathrm{off}}(f)\Delta t],
 \quad P_{\mathrm{on}}=1-\exp(-k_{\mathrm{on}}\Delta t).
 \label{eq:bell}
\end{equation}
The independent hypothesis uses
\begin{equation}
 F_{\mathrm{site}}=\sum_{j\in\mathrm{bound}}k_cx_j,
 \label{eq:independent}
\end{equation}
whereas the equal-load-sharing hypothesis uses
\begin{equation}
 f_i=\frac{F_{\mathrm{site}}}{i},\qquad
 r_i=i k_{\mathrm{off}}^0\exp\!\left(\frac{F_{\mathrm{site}}}{iF_b}\right).
 \label{eq:shared}
\end{equation}
An empty angular sector cannot bind or load a clutch. This active-mask guard prevents phantom
traction against nonexistent collagen. Bell rupture of a cell--fibre clutch produces adhesion
slippage; Bell rupture of a fibre--fibre crosslink in G5E changes ECM topology. They are separate
processes with separate parameters.

\section{G5D: collective invasion versus single-cell escape}

Every rigid cell disk moves under its cell--ECM reaction force and cell--cell mechanics:
\keylabel
\begin{equation}
 \zeta_c\dot{\cv}_a=
 \Fv_a^{\mathrm{reaction}}+
 \sum_{b\ne a}\Fv_{a\leftarrow b}^{cc},
 \qquad
 \Fv_a^{\mathrm{reaction}}=-\sum_{s\in a}\Fv_s^{\mathrm{cell\to ECM}}.
 \label{eq:cellmotion}
\end{equation}
A cell reels collagen inward; the opposite reaction tends to pull the cell toward the gripped
ECM. Adhesion resists separation. The velocity is capped when necessary:
\begin{equation}
 \mathbf v_a=\Fv_a^{\mathrm{net}}/\zeta_c,\qquad
 \lVert\mathbf v_a\rVert>v_{\max}\Rightarrow
 \mathbf v_a\leftarrow v_{\max}\frac{\mathbf v_a}{\lVert\mathbf v_a\rVert}.
 \label{eq:speedcap}
\end{equation}
High adhesion favours cohesion; low adhesion permits cell--ECM reaction to separate individual
cells more easily. ``Invasion'' in this model means outward displacement arising from this
force competition, not a complete biological tissue-invasion programme.

Optional leader heterogeneity is imposed on selected outer cells:
\begin{equation}
 s_a=s_{\mathrm{leader}}<1,\qquad
 k_{\mathrm{adh},ab}=k_{\mathrm{adh}}\min(s_a,s_b),\qquad
 F_a^{\mathrm{pull}}=\lambda_aF^{\mathrm{pull}}.
 \label{eq:leader}
\end{equation}
This is a modelling hypothesis; the current outermost-fraction selection is deterministic and is
not an emergent EMT mechanism.

\section{G5E: crosslink topology plasticity}

Loaded fibre--fibre crosslinks may rupture with a Bell-type hazard:
\keylabel
\begin{equation}
 k_{\mathrm{off},x}(F_x)=k_{\mathrm{off},x}^0
 \exp\!\left(\frac{F_x}{F_{b,x}}\right),
 \qquad
 P_{\mathrm{rupture}}=1-
 \exp[-k_{\mathrm{off},x}(F_x)\Delta t_x].
 \label{eq:xrupture}
\end{equation}
Current defaults $k_{\mathrm{off},x}^0=2\times10^{-4}\unit{s^{-1}}$ and
$F_{b,x}=3\unit{nN}$ are explicit assumptions pending calibration. When a new crossing is
accepted, its current configuration becomes stress free:
\begin{equation}
 \mathbf d_{x,0}\leftarrow\mathbf p_b(t)-\mathbf p_a(t),
 \qquad\text{with re-weld acceptance probability }p_{\mathrm{reform}}.
 \label{eq:reweld}
\end{equation}
Changing which fibres are welded and where the new rest state is created is topology-changing
remodelling, not SLS stress relaxation.

Topology memory and residual response are summarized by
\begin{align}
 \kappa_{\mathrm{topo}}&=1-
 \frac{|L_0\cap L_T|}{|L_0|},\\
 \kappa_{\mathrm{rms}}&=u_{\mathrm{residual}}/u_{\mathrm{peak}},\\
 \kappa_{\mathrm{order}}&=
 \frac{S_{\mathrm{final}}-S_0}{S_{\mathrm{peak}}-S_0}.
 \label{eq:memory}
\end{align}
$\kappa_{\mathrm{topo}}$ is the cleanest topology metric. Residual displacement alone can be
confounded by slow elastic relaxation and near-zero modes.

\section{Ablation logic}

An ablation is a matched comparison in which one mechanism is removed or activated while the
initial network, seed, duration, and other parameters remain fixed:
\keylabel
\begin{equation}
 \Delta Y_{(-m)}=Y_{\mathrm{baseline}}-Y_{\mathrm{mechanism}\ m\ \mathrm{removed}}.
 \label{eq:ablation}
\end{equation}
$Y$ may be invasion distance, cohesion, radial order, displacement reach, traction, or topology
memory. The comparison attributes effects inside the model; it does not by itself establish an
experimental effect size.

\begin{longtable}{>{\raggedright\arraybackslash}p{0.22\textwidth}p{0.33\textwidth}p{0.35\textwidth}}
\toprule
Ablation & Mechanism isolated & Primary readout \\
\midrule
No clutch & Active cell--fibre loading & Invasion, traction, and slip events. \\
No adhesion & Organoid cohesion & Dispersion, escaped-cell fraction, and collective integrity. \\
No crosslinks & Network force transmission & Indirect displacement, radial order, and invasion. \\
Stiffening off/on & Nonlinear tensile recruitment & Alignment and displacement reach. \\
Plasticity off/on & Topology memory & $\kappa_{\mathrm{topo}}$ and post-unload persistence. \\
Leaders off/on & Prescribed cell heterogeneity & Leader displacement and follower cohesion. \\
\bottomrule
\end{longtable}

In the current matched two-hour runs, approximate reported displacement/invasion is
$25\unit{\mu m}$ for the clutch baseline, $4.7\unit{\mu m}$ with no clutch,
$68\unit{\mu m}$ with no adhesion but worse cohesion, $7.7\unit{\mu m}$ with no crosslinks,
$24\unit{\mu m}$ with strain stiffening, and $25\unit{\mu m}$ with plasticity. These are
single-model outputs, not experimental effect sizes; ensembles and calibration are required.

\section{Parameter map and evidence boundaries}

\begin{longtable}{>{\raggedright\arraybackslash}p{0.25\textwidth}p{0.24\textwidth}p{0.41\textwidth}}
\toprule
Parameter & Current default & Role and status \\
\midrule
Cell radius $R_c$ & $9\unit{\mu m}$ & Single-cell scale, not a fit to one cell line. \\
Organoid centre radius & $63\unit{\mu m}$ & Two-dimensional core geometry. \\
Cell spacing $d_0$ & $18\unit{\mu m}$ & Touching hexagonal disks and zero pair force. \\
Cell repulsion / adhesion & 40 / $6\unit{nN/\mu m}$ & Effective non-overlap and cohesion parameters. \\
Adhesion range & $6\unit{\mu m}$ & Finite-neighbour interaction assumption. \\
Collagen modulus & $3\unit{MPa}$ & Inherited softened mechanism-first choice. \\
Crosslink stiffness & $10\unit{nN/\mu m}$ & Effective link penalty. \\
Fibres / domain & 420 / $440\unit{\mu m}$ & Organoid-scale computational network. \\
Crosslink retain fraction & 0.85 & Generator connectivity target, not G4's probability mechanism. \\
Pull per cell & $12\unit{nN}$ & Constant-pull branch only. \\
Bead time step & $0.05\unit{s}$ & Inherited numerical resolution. \\
Cell drag / speed cap & $400\unit{nN\,s/\mu m}$ / $0.02\unit{\mu m/s}$ & Effective mobility and numerical/biological guard. \\
\bottomrule
\end{longtable}

The inherited effective clutch defaults are $N=12$, $k_c=2\unit{nN/\mu m}$,
$k_{\mathrm{on}}=0.055\unit{s^{-1}}$, $k_{\mathrm{off}}^0=0.018\unit{s^{-1}}$,
$F_b=1.5\unit{nN}$, $v_0=0.025\unit{\mu m/s}$, and
$F_{\mathrm{stall}}=8\unit{nN/site}$. G4 crosslink probability and G5 crosslink retention
fraction are different constructions and should not be compared as one biological density.

G5 remains a 2-D rigid-disk model without cell deformation, a nucleus, division, death, growth
pressure, proteolysis, nutrient fields, three-dimensional pore confinement, or an emergent
organoid polarity field. Several coefficients are exploratory. Matched ensembles and sensitivity
analysis are required before biological prediction.

\subsection*{Key references}
\begin{enumerate}[leftmargin=1.5em,itemsep=2pt]
\item Lee et al. (2014), \emph{PLOS ONE} 9:e111896. doi:10.1371/journal.pone.0111896.
\item Abhilash et al. (2014), \emph{Biophysical Journal} 107:1829--1840. doi:10.1016/j.bpj.2014.08.029.
\item Wang et al. (2014), \emph{Biophysical Journal} 107:2592--2603. doi:10.1016/j.bpj.2014.09.044.
\item Bell (1978), \emph{Science} 200:618--627. doi:10.1126/science.347575.
\item Bangasser and Odde (2013), \emph{Cellular and Molecular Bioengineering} 6:449--459.
\item Erdmann and Schwarz (2004), arXiv:cond-mat/0403552.
\item Ilina and Friedl (2020), \emph{Nature Cell Biology}. doi:10.1038/s41556-020-0552-6.
\item Steinwachs et al. (2016), \emph{Nature Methods} 13:171--176. doi:10.1038/nmeth.3685.
\item Mark et al. (2020), \emph{eLife} 9:e51912.
\item Nam et al. (2016), \emph{Biophysical Journal} 111:2296--2308.
\item Wisdom et al. (2018), \emph{Nature Communications}. doi:10.1038/s41467-018-06641-z.
\item Kim/Ahn et al. (2019), \emph{Acta Biomaterialia} 84:280. PMID:30500449.
\item Carey et al. (2016), \emph{Integrative Biology} 8:821--835. doi:10.1039/C6IB00030D.
\end{enumerate}

\medskip
\noindent\textbf{Claim boundary.} G5 is a controlled 2-D mechanism platform showing how
inherited clutch traction, cell--cell cohesion, and crosslinked collagen jointly produce
different organoid behaviours. It is not yet a predictive digital twin of tumour invasion.

\end{document}
